\documentclass[11pt]{article}
\usepackage{jtsguide}
\setshorttitle{TestBuilder \& TestRunner}

\title{JTS Topology Suite --- TestRunner \& TestBuilder User Guide}
\author{Jonathan Aquino (2001--2005) \\ rebuilt from the 2005 text}
\date{}

\begin{document}
\jtscover{TestRunner \& TestBuilder User Guide}{Rebuilt from the 3 February 2005 text}{Base text: 3 February 2005 \textperiodcentered\ rebuild: 16 August 2026}

\noindent
This guide describes how to use the JTS TestRunner, a Java application
that validates the JTS Topology Suite, and the JTS TestBuilder, a Java
application that creates, visualizes, and runs geometry tests. The 2005
chapter plan is kept. Startup, packages, and curve drawing are brought
up to this tree.

Package names in this tree are \texttt{org.locationtech.jts}. The 2005
copies used \texttt{com.vividsolutions.jts}. LocationTech JTS remains
the upstream project; this fork is not described as if it were already
upstream.

\clearpage
\changectrl{%
  1.0 & 13 Sep 2001 & Jonathan Aquino & Original draft. \\
  1.1 & 17 Dec 2001 & Jonathan Aquino & Added TestBuilder documentation. \\
  1.2 & 30 Jan 2002 & Jonathan Aquino & Added ACME licence (historical). \\
  1.2-rebuild & 16 Aug 2026 & This tree & LaTeX rebuild. Startup, packages,
    curve drawing, WKT apply, A/B colours, OverlayNGCurve, Hausdorff.
    ACME appendix replaced with the JTS dual licence. \\
}

\tableofcontents
\clearpage

\section{Overview}
The TestRunner and TestBuilder are used by people who are developing
or evaluating JTS. While similarly named, the JTS TestRunner is not
the JUnit TestRunner. JUnit is a general Java testing framework. The
JTS TestRunner is easier for testing JTS computations because tests
are specified in XML text files --- no Java coding or compilation.
Those files can be created by hand or generated by the TestBuilder.

The 2005 guide required Java 2 JRE/SDK v1.3. This tree targets Java 8
and above. Download a current JDK from the vendor of your choice; the
2005 \texttt{java.sun.com/products} link is historical.

The 2005 class names were
\texttt{com.vividsolutions.jtstest.testrunner.TopologyTestApp} and
\texttt{com.vividsolutions.jtstest.testbuilder.JTSTest}. On this tree
they are
\texttt{org.locationtech.jtstest.testrunner.JTSTestRunnerCmd} and
\texttt{org.locationtech.jtstest.testbuilder.JTSTestBuilder}. The
usual way to start TestBuilder is the built jar, not a hand-built
classpath of \texttt{jdom.jar} / \texttt{xerces.jar} / \texttt{jts.jar}
/ \texttt{jts\_test.jar}.

This guide is divided into two main sections: Using the TestRunner,
and Using the TestBuilder. An appendix details the XML test-file
format. For more information on JTS, see the JTS Technical
Specifications.

\section{Using the TestRunner}
This section describes how to use the TestRunner to run a test file
that validates the JTS Topology Suite. Test files are created using
the TestBuilder (Section~\ref{sec:tb}) or by hand
(Section~\ref{sec:xml}).

\subsection{Starting}
From a built tree, run the TestRunner via the scripts in
\texttt{bin/} or:

\begin{lstlisting}[language=bash]
java -cp modules/tests/target/jts-tests-*-SNAPSHOT.jar \
     org.locationtech.jtstest.testrunner.JTSTestRunnerCmd \
     -files path/to/tests.xml
\end{lstlisting}

The 2005 GUI mode (\texttt{-gui} / \texttt{testrunner.bat}) and the
switches \texttt{-files}, \texttt{-properties}, \texttt{-verbose} are
the same idea: a list of XML files, a results panel or text listing,
and a status line of cases / tests / passed / failed / exceptions.

Graphical mode has file-chooser dialogs that make it easy to build a
list of test files. Text mode is quicker to start. The
\texttt{-files} switch can be used to run all test files in a given
directory.

\subsection{Reading the results}
A case specifies two geometries (A and B). A test specifies an
operation and its expected value. The status line reports counts of
cases, tests, passes, failures, thrown exceptions, and parse
exceptions.

\begin{center}
\begin{tabular}{>{\raggedright\arraybackslash}p{3.4cm}p{9.8cm}}
  \toprule
  \textbf{Statistic} & \textbf{Meaning} \\
  \midrule
  Cases            & Each case specifies two geometries to test. \\
  Tests            & Each test specifies an operation on those
                     geometries. \\
  Passed           & Actual value matches expected value. \\
  Failed           & Actual value differs from expected value. \\
  Exception        & An error occurred inside JTS or the TestRunner. \\
  Parse exception  & The test file has a syntax error. \\
  \bottomrule
\end{tabular}
\end{center}

Passed-test detail is shown only with \texttt{-verbose}. A failed
test means either fix the test or fix JTS. A parse exception is a
problem in the XML or WKT (a missing comma in a LINESTRING is the
2005 example and is still the usual cause).

\subsection{Switches}
\begin{center}
\begin{tabular}{>{\raggedright\arraybackslash}p{3.0cm}p{10.2cm}}
  \toprule
  \textbf{Switch} & \textbf{Description} \\
  \midrule
  \texttt{-files}      & List of test files and directories. Separate
                         items with spaces. \\
  \texttt{-properties} & A \texttt{.properties} file that remembers
                         the last file list. \\
  \texttt{-gui}        & Graphical mode, when the build still ships
                         that entry point. \\
  \texttt{-verbose}    & Show passed tests as well as failures. \\
  \bottomrule
\end{tabular}
\end{center}

The easiest way to specify every file in a directory is to pass the
directory to \texttt{-files}. Any combination of switches is
permitted. If no switches are specified, a description of the
switches is displayed.

\begin{amendment}
XML \texttt{<a>} / \texttt{<b>} bodies that contain
\texttt{CIRCULARSTRING} / \texttt{COMPOUNDCURVE} /
\texttt{CURVEPOLYGON} / \texttt{MULTICURVE} / \texttt{MULTISURFACE}
must be parsed with a curve-capable reader. On this tree the
TestRunner / TestCase path uses \texttt{CurveWKTReader}.

A test that calls \texttt{Geometry.intersection} on a curve is not
\texttt{OverlayNGCurve}. Name the OverlayNGCurve function if that is
what you mean to assert.

A \texttt{DiscreteHausdorffDistance} assertion is the public API: two
certified pairs are closed-form; everything else is chords.
\texttt{fail()} stays on the general case.
\end{amendment}

\section{Using the TestBuilder}
\label{sec:tb}
The 2005 guide called TestBuilder experimental, with a useful starting
point for creating, visualizing, and running test files. That
description is still fair. The application on this tree is the
current LocationTech JTS TestBuilder plus the curve tools that are
actually merged.

Useful features:
\begin{itemize}
  \item creating test geometries visually or by entering Well-Known
        Text;
  \item snapping to a grid;
  \item specifying expected values for spatial functions and the
        intersection matrix;
  \item opening and saving XML test files.
\end{itemize}

The 2005 ``Save As HTML'' bulk documentation generator is not claimed
here as a supported workflow. The \texttt{vividsolutions.com/jts/tests}
URL in the 2005 copy is historical.

\subsection{Starting the TestBuilder}
Build the app module, then from the project root:

\begin{lstlisting}[language=bash]
java -jar modules/app/target/JTSTestBuilder.jar -Xmx2000M
\end{lstlisting}

On macOS, if the native look-and-feel misbehaves:

\begin{lstlisting}[language=bash]
java -Dswing.defaultlaf=javax.swing.plaf.metal.MetalLookAndFeel \
     -jar modules/app/target/JTSTestBuilder.jar -Xmx2000M
\end{lstlisting}

Do not start it as the 2005
\texttt{java com.vividsolutions.jtstest.testbuilder.JTSTest} with a
four-jar classpath. That layout is gone.

\subsection{Opening a test file}
File / Open XML File(s). The TestBuilder displays one test case at a
time: a pair of geometries and their tests. Use Previous / Next or
the Tests list to move. Pan, zoom, zoom 1:1, and zoom-to-extent
examine the geometries. The WKT tab shows Well-Known Text.

\subsection{Precision model}
A test file's precision model is set from the TestBuilder
precision-model dialog (Edit / Precision Model). The default is
Floating. See the Technical Specifications.

\subsection{Creating geometries}
The 2005 workflow is still the right picture: pick geometry A or B,
pick a draw tool, click in the view. Geometry A is shown in blue;
geometry B, in red. You can also type WKT. Erase, undo point, and
undo part still apply.

On this tree the draw tools include the linear set (point, linestring,
polygon, rectangle) and the curve set that is actually merged:
\texttt{CircularString}, \texttt{CompoundCurve},
\texttt{CurvePolygon}. Triangle and TIN tools also exist.

Draw-tool icons follow the A/B focus: index 0 (A) is blue-dominant;
otherwise (B) is red-dominant. That applies to the linear draw tools
and to CircularString / CompoundCurve / CurvePolygon.

WKT apply uses \texttt{CurveWKTReader}, so
\texttt{CIRCULARSTRING} pasted into the WKT tab is a
\texttt{CircularString}, not a parse error. Enter in the A or B WKT
pane loads that pane. The Clear control is labeled and does not wipe
the other pane on apply.

\fighole{UG-1 draw CS}
\fighole{UG-2 A/B colors}
\fighole{UG-3 CompoundCurve}
\fighole{UG-4 CurvePolygon disc}

A clothoid may be added as a non-leading
\texttt{COMPOUNDCURVE} member from the CompoundCurve tool and from
the clothoid playground extras. That path is extras, not a
closed-form construction.

\subsection{Predicates and functions}
The 2005 Predicates tab (expected DE-9IM, binary predicates) and
Functions tab (union, intersection, buffer, \ldots) are the same
idea. The current UI is a function tree: pick a function, run it,
inspect the result geometry and its WKT.

\texttt{OverlayNGCurve} is registered as its own function class
(intersection / union / difference / symDifference). That is the
curve overlay surface. \texttt{Geometry.intersection} in the Geometry
function group is not OverlayNGCurve.

Closed-form constructions exposed as functions include disc and
single-arc convex hull, the compound-curve hull of
circular-plus-straight members, \texttt{MaximumInscribedCircle} on a
certified disc and on a stadium, and \texttt{LargestEmptyCircle} on
legal curve obstacles.

\subsection{Saving}
File / Save As XML writes the current cases. Default / export WKB
paths route through \texttt{CurveWKBWriter}, so a saved curve keeps
ISO codes 8--12. A bare \texttt{WKBWriter.write} call still emits
type 2 or 3.

\section{Curves in TestBuilder}
TestBuilder can draw \texttt{CircularString}, \texttt{CompoundCurve},
and \texttt{CurvePolygon} (toolbar tools plus WKT via
\texttt{CurveWKTReader}). It also exposes OverlayNGCurve and the
distance functions.

On this tree the concrete SFA/SQL-MM curve types live in the
\texttt{jts-curve} module
(\texttt{org.locationtech.jts.geom.curve}). Construction is opt-in:
the default \texttt{GeometryFactory} throws on the curve
\texttt{create*} methods. Use \texttt{CurveGeometryFactory}.
Constructors do not linearise.

Core \texttt{WKTReader} still rejects \texttt{CIRCULARSTRING} and the
other SQL-MM keywords. Read them with \texttt{CurveWKTReader}. Write
structured members with \texttt{CurveWKTWriter}.

WKB type codes 8--12 are first-class in core \texttt{WKBReader}.
\texttt{CurveWKBWriter} emits codes 8--12. A bare \texttt{WKBWriter}
still emits type 2 (CircularString / CompoundCurve) or 3
(CurvePolygon).

The curve overlay type is \texttt{OverlayNGCurve}. There is no public
noder. Closed-form overlay answers are limited to certified discs.

Public \texttt{DiscreteHausdorffDistance} (commit \texttt{0ca71b})
has a closed form for two pairs only: (1)~a single-arc
\texttt{CircularString} toward a single-segment \texttt{LineString},
apex $\sqrt{949}/6 - 7/6 \approx 3.967641$ for
\texttt{CIRCULARSTRING (0 0, 2 3, 10 0)} vs
\texttt{LINESTRING (0 0, 10 0)}; (2)~two circular discs. A
\texttt{MultiSurface} unwrap of one disc is the same pair. The exact
path skips densify; densify is not the laser. For every other pair
the public API still sees vertices / chords. Fr\'echet remains open.

\fighole{UG-5 APEX}
\fighole{UG-6 two discs}
\fighole{UG-7 disc + single-arc hull}

\section{Appendix: test file format}
\label{sec:xml}
A test file is XML. It contains one or more cases; each case contains
one or more tests. A case specifies two geometries. A test specifies
an operation and its expected result.

\subsection{Example}
\begin{lstlisting}[language=XML]
<run>
  <desc>example</desc>
  <precisionModel type="FLOATING"/>
  <case>
    <desc>point in a polygon</desc>
    <a>POLYGON ((0 0, 10 0, 10 10, 0 10, 0 0))</a>
    <b>POINT (5 5)</b>
    <test>
       <op name="contains">true</op>
    </test>
  </case>
  <case>
    <desc>single-arc D-HF witness</desc>
    <a>CIRCULARSTRING (0 0, 2 3, 10 0)</a>
    <b>LINESTRING (0 0, 10 0)</b>
    <test>
       <op name="orientedDistance" arg1="a" arg2="b">3.967641</op>
    </test>
  </case>
</run>
\end{lstlisting}

The second case is the D-HF closed-form witness on this tree (apex
$\sqrt{949}/6 - 7/6$). A case that is not one of the two certified
pairs must not be written as if
\texttt{DiscreteHausdorffDistance} were arc-length Hausdorff. The
first case is the 2005 example, unchanged.

\subsection{XML shape}
The 2005 DTD-style listing is still the right shape: \texttt{<run>}
contains \texttt{<precisionModel>} and one or more \texttt{<case>};
each case has \texttt{<a>}, optional \texttt{<b>}, and
\texttt{<test>} / \texttt{<op>} with a name and expected value. WKT
inside \texttt{<a>} / \texttt{<b>} may now be a curve keyword; the
reader on this tree is \texttt{CurveWKTReader}.

\section{Appendix: licence}
JTS is dual-licensed under the Eclipse Public License 2.0 and the
Eclipse Distribution License 1.0. See \texttt{LICENSE\_EPLv2.txt},
\texttt{LICENSE\_EDLv1.txt}, and \texttt{LICENSES.md}. The 2005
``Appendix: ACME Licence'' is not reproduced. It was never the JTS
project licence on this tree.

\vspace{1.2em}
\noindent\textit{End of TestRunner \& TestBuilder User Guide. Named
figure holes are empty 4:3 slots (target 1600$\times$1200, full
TestBuilder window). 2005 Amyuni screen-shots are omitted; the
procedures they illustrated are kept in prose.}

\end{document}
